% ===========================================================================
%  Jàng Afrig — CM1 : Le Web et HTML
%  Licence 1 DUT1-Informatique — ESP/UCAD — 2025–2026
%  Dr. El Hadji Bassirou TOURÉ
%  Compilé avec : xelatex (2–3 passes)
% ===========================================================================
\documentclass[11pt,a4paper]{article}

\usepackage{fontspec}
\setmainfont{Latin Modern Roman}
\setsansfont{Latin Modern Sans}
\setmonofont{Latin Modern Mono}
\usepackage{polyglossia}
\setdefaultlanguage{french}

\usepackage[margin=2.5cm]{geometry}
\usepackage{amsmath,amssymb}
\usepackage{graphicx}
\graphicspath{{figures/}}
\usepackage{xcolor}
\usepackage{booktabs}
\usepackage{hyperref}
\usepackage{tcolorbox}
\tcbuselibrary{theorems,breakable,skins}
\usepackage{enumitem}
\usepackage{fancyhdr}
\usepackage{array}
\usepackage{pifont}
\usepackage{longtable}
\usepackage{multirow}

% --- TikZ pour les schémas ---
\usepackage{tikz}
\usetikzlibrary{shapes.geometric,arrows.meta,positioning,calc,backgrounds,fit,shadows.blur,trees}

\newcommand{\cmark}{\ding{51}}
\newcommand{\xmark}{\ding{55}}

\hypersetup{
    colorlinks=true,
    linkcolor=blue!60!black,
    urlcolor=blue!60!black,
    citecolor=blue!60!black
}

\newtcolorbox{objectifbox}[1][]{colback=blue!5,colframe=blue!50,fonttitle=\bfseries,title=#1,breakable}
\newtcolorbox{infobox}[1][]{colback=green!5,colframe=green!50!black,fonttitle=\bfseries,title=#1,breakable}
\newtcolorbox{attentionbox}[1][]{colback=orange!5,colframe=orange!70,fonttitle=\bfseries,title=#1,breakable}
\newtcolorbox{retenirbox}[1][]{colback=yellow!10,colframe=yellow!50!black,fonttitle=\bfseries,title=#1,breakable}
\newtcolorbox{prereqbox}[1][]{colback=cyan!5,colframe=cyan!50!black,fonttitle=\bfseries,title=#1,breakable}
\newtcolorbox{projetbox}[1][]{colback=teal!5,colframe=teal!50!black,fonttitle=\bfseries,title=#1,breakable}

% --- Couleurs pour les schémas ---
\definecolor{schclient}{RGB}{45,130,200}
\definecolor{schserveur}{RGB}{200,80,60}
\definecolor{schhtml}{RGB}{228,77,38}
\definecolor{schhead}{RGB}{110,70,180}
\definecolor{schbody}{RGB}{45,160,110}

% --- Code source ---
\usepackage{listings}
\lstdefinelanguage{HTML5}{
  morekeywords={<!DOCTYPE,html,head,meta,title,body,h1,h2,h3,p,ul,ol,li,table,tr,th,td,a,img,strong,em,br,div,span},
  sensitive=true,
  morecomment=[s]{<!--}{-->},
  morestring=[b]",
}
\lstset{
  basicstyle=\ttfamily\small,
  backgroundcolor=\color{gray!8},
  frame=single,
  framerule=0.5pt,
  rulecolor=\color{gray!40},
  breaklines=true,
  breakatwhitespace=true,
  showstringspaces=false,
  tabsize=2,
  xleftmargin=4pt,
  xrightmargin=4pt,
  aboveskip=8pt,
  belowskip=8pt,
  literate={é}{{\'e}}1 {è}{{\`e}}1 {ê}{{\^e}}1 {à}{{\`a}}1 {ù}{{\`u}}1
           {ô}{{\^o}}1 {î}{{\^i}}1 {ç}{{\c{c}}}1 {É}{{\'E}}1
           {→}{{$\rightarrow$}}1
}

\pagestyle{fancy}
\fancyhf{}
\fancyhead[L]{\small Jàng Afrig --- CM1 : Le Web et HTML}
\fancyhead[R]{\small ESP/UCAD --- L1 DUT1-Info}
\fancyfoot[C]{\thepage}
\fancyfoot[L]{\small Dr. El Hadji Bassirou TOURÉ}
\fancyfoot[R]{\small 2025--2026}
\renewcommand{\headrulewidth}{0.4pt}
\renewcommand{\footrulewidth}{0.4pt}

\title{%
\vspace{-1cm}
\textbf{\LARGE Fondamentaux du Web}\\
\textbf{\LARGE HTML, CSS, JavaScript}\\[0.8cm]
\textbf{\huge CM1 --- Le Web et HTML}\\[0.3cm]
\textbf{\Large Comment fonctionne une page web, et comment on la structure}\\[0.5cm]
\large Licence 1 DUT1-Informatique --- 2025--2026}
\author{Dr. El Hadji Bassirou TOURÉ\\
École Supérieure Polytechnique\\
Département Génie Informatique\\
Université Cheikh Anta Diop de Dakar}
\date{}

\begin{document}

\maketitle
\thispagestyle{empty}

\vspace{0.5cm}

% ===========================================================================
% BOÎTE PRÉREQUIS
% ===========================================================================
\begin{prereqbox}[Informations pratiques]
\begin{center}
\begin{tabular}{ll}
\toprule
\textbf{Type} & Cours Magistral (CM) \\
\textbf{Séance} & S1 --- première séance du semestre \\
\textbf{Prérequis} & Lab 0 terminé (environnement installé + repo \texttt{jang-afrig}) \\
\textbf{Ce qu'on verra après} & Lab 1 : Ma première page HTML (S2) \\
\textbf{Devoir} & Aucun. Juste lire le PDF du Lab 1 pour la prochaine séance. \\
\bottomrule
\end{tabular}
\end{center}
\end{prereqbox}

\begin{objectifbox}[Objectifs du CM]
À la fin de cette séance, vous saurez :
\begin{enumerate}[nosep]
    \item Expliquer \textbf{comment une page web arrive} dans votre navigateur
    \item Distinguer \textbf{les trois rôles} : HTML, CSS, JavaScript
    \item Décrire \textbf{l'anatomie} d'une page HTML (DOCTYPE, \texttt{head}, \texttt{body})
    \item Utiliser les \textbf{balises de contenu} essentielles (titres, paragraphes, listes, tableaux, liens, images)
    \item Imbriquer des listes et fusionner des cellules de tableau (\texttt{colspan}, \texttt{rowspan})
    \item Comprendre \textbf{ce qu'on va construire} : Jàng Afrig
\end{enumerate}
\end{objectifbox}

\tableofcontents

\newpage

% ===========================================================================
\section{Une page web, c'est quoi ?}
% ===========================================================================

Ouvrez un navigateur, tapez \url{https://ucad.sn} et appuyez sur Entrée.
Vous voyez une page avec des titres, du texte, des images, des menus, des
couleurs, des animations. Cette page n'existe pas par magie : quelqu'un
l'a \textbf{construite}. Et bientôt, vous saurez la construire aussi.

\subsection{L'analogie de la maison}

Une page web est comme une maison. Trois technologies la composent, chacune
avec un rôle précis :

\begin{center}
\begin{tabular}{@{}lll@{}}
\toprule
\textbf{Technologie} & \textbf{Rôle} & \textbf{Dans la maison} \\
\midrule
\textbf{HTML} & Structure & Murs, cloisons, fondations \\
\textbf{CSS}  & Décoration & Peinture, meubles, couleurs \\
\textbf{JavaScript} & Comportement & Électricité, eau, interrupteurs \\
\bottomrule
\end{tabular}
\end{center}

\medskip

\begin{itemize}[leftmargin=*,nosep]
    \item Une page \textbf{HTML seule} est une maison en briques brutes : fonctionnelle mais austère.
    \item Une page \textbf{HTML + CSS} est une maison décorée : on aime y être.
    \item Une page \textbf{HTML + CSS + JavaScript} est une maison \textit{vivante} : on peut allumer la lumière, monter le son, ouvrir l'eau.
\end{itemize}

\begin{retenirbox}[À retenir]
\textbf{HTML structure} le contenu. \textbf{CSS} l'habille. \textbf{JavaScript} le rend interactif.

Dans ce cours, on apprend les trois, \textbf{dans cet ordre}. Aujourd'hui, on pose les fondations : \textbf{HTML}.
\end{retenirbox}

\subsection{Pourquoi commencer par HTML ?}

Parce qu'on ne décore pas une maison qu'on n'a pas encore construite. HTML
est le \textbf{socle}. Sans HTML, ni CSS ni JavaScript ne peuvent exister.
Les trois semaines qui suivent porteront intégralement sur HTML. Ensuite
CSS, ensuite JavaScript.

\newpage

% ===========================================================================
\section{Comment une page web arrive dans mon navigateur ?}
% ===========================================================================

Vous tapez \url{https://ucad.sn} dans votre navigateur. En moins d'une seconde,
la page s'affiche. Que s'est-il passé ? Votre navigateur a discuté avec un
\textbf{serveur} quelque part dans le monde.

\subsection{Le modèle client / serveur}

\begin{figure}[!ht]
\centering
\begin{tikzpicture}[
    node distance=2cm,
    acteur/.style={rounded corners=8pt, draw, thick, minimum width=4cm, minimum height=2.2cm, align=center, font=\bfseries},
    fleche/.style={-{Stealth[length=3.5mm]}, thick, line width=1.2pt},
    lib/.style={midway, font=\small\itshape, fill=white, inner sep=3pt}
]
    \node[acteur, fill=schclient!15, draw=schclient] (client)
      {Votre navigateur\\[2pt]\normalfont\small (Chrome, Firefox)\\ \normalfont\scriptsize à Dakar, sur votre PC};

    \node[acteur, fill=schserveur!15, draw=schserveur, right=5cm of client] (serveur)
      {Le serveur\\[2pt]\normalfont\small (un ordinateur)\\ \normalfont\scriptsize quelque part dans le monde};

    % Flèche aller (requête)
    \draw[fleche, color=schclient]
      ([yshift=0.5cm]client.east) -- node[lib, above] {1. Donne-moi la page ucad.sn}
      ([yshift=0.5cm]serveur.west);

    % Flèche retour (réponse)
    \draw[fleche, color=schserveur]
      ([yshift=-0.5cm]serveur.west) -- node[lib, below] {2. Voici le code HTML}
      ([yshift=-0.5cm]client.east);

\end{tikzpicture}
\caption{Le dialogue entre votre navigateur (client) et un serveur.}
\label{fig:client-serveur}
\end{figure}

\begin{enumerate}[leftmargin=*,nosep]
    \item Vous tapez une adresse. Votre navigateur envoie une \textbf{requête} au serveur qui héberge le site.
    \item Le serveur répond en vous envoyant du \textbf{code HTML} (plus parfois du CSS, du JavaScript, des images).
    \item Votre navigateur \textbf{lit} ce code HTML et \textbf{affiche} une page à l'écran.
\end{enumerate}

\begin{infobox}[L'analogie du restaurant]
Vous entrez dans un restaurant (\textbf{le navigateur ouvre le site}).
Vous dites au serveur : «\,Je voudrais le thiéboudienne\,» (\textbf{la requête}).
Le serveur va en cuisine, prépare votre plat et vous l'apporte (\textbf{la réponse du serveur web}).
Vous mangez ce qu'on vous a servi (\textbf{le navigateur affiche la page}).
\end{infobox}

\subsection{Ce que le serveur envoie : du texte, rien que du texte}

C'est le point étonnant : quand le serveur répond, il envoie juste un
\textbf{gros fichier texte}. Ce fichier s'appelle du \textbf{HTML}. C'est
ensuite votre navigateur qui transforme ce texte en page visuelle.

\begin{retenirbox}[À retenir]
Le serveur envoie du \textbf{texte} (HTML). Le navigateur transforme ce
texte en \textbf{page visuelle}. Votre travail de développeur web :
\textbf{écrire ce texte}.
\end{retenirbox}

\newpage

% ===========================================================================
\section{Anatomie d'une page HTML}
% ===========================================================================

Une page HTML a toujours la même forme. Comme une lettre officielle a un
en-tête, un corps et une signature, une page HTML a trois parties :

\begin{figure}[!ht]
\centering
\begin{tikzpicture}[
    level distance=1.7cm,
    sibling distance=5cm,
    every node/.style={rounded corners=4pt, draw, thick, font=\ttfamily, inner sep=6pt},
    edge from parent/.style={draw, thick, -{Stealth[length=2mm]}}
]
    \node[fill=schhtml!15, draw=schhtml, text=schhtml] {\textbf{<html>}}
        child {node[fill=schhead!15, draw=schhead, text=schhead] {\textbf{<head>}}
            child[sibling distance=2.2cm] {node[fill=white, draw=schhead!50] {\small <title>}}
            child[sibling distance=2.2cm] {node[fill=white, draw=schhead!50] {\small <meta>}}
        }
        child {node[fill=schbody!15, draw=schbody, text=schbody] {\textbf{<body>}}
            child[sibling distance=1.8cm] {node[fill=white, draw=schbody!50] {\small <h1>}}
            child[sibling distance=1.8cm] {node[fill=white, draw=schbody!50] {\small <p>}}
            child[sibling distance=1.8cm] {node[fill=white, draw=schbody!50] {\small <img>}}
        };
\end{tikzpicture}
\caption{La structure d'une page HTML : deux grandes zones, \texttt{head} et \texttt{body}.}
\label{fig:anatomie}
\end{figure}

\subsection{Le squelette minimal}

Voici le \textbf{squelette minimal} de toute page HTML. Cette structure se
retrouve sur \emph{chaque} site web que vous avez ouvert dans votre vie.

\begin{lstlisting}[language=HTML5]
<!DOCTYPE html>
<html lang="fr">
    <head>
        <meta charset="UTF-8">
        <title>Ma page</title>
    </head>
    <body>
        <h1>Bonjour !</h1>
        <p>Ceci est ma première page web.</p>
    </body>
</html>
\end{lstlisting}

\subsection{Ce que fait chaque ligne}

\begin{center}
\begin{tabular}{@{}lp{9cm}@{}}
\toprule
\textbf{Ligne} & \textbf{Rôle} \\
\midrule
\texttt{<!DOCTYPE html>} & Dit au navigateur : «\,Je suis une page HTML moderne (HTML5).\,» \\
\addlinespace
\texttt{<html lang="fr">} & Démarre la page. \texttt{lang="fr"} indique la langue (utile pour les moteurs de recherche et les lecteurs d'écran). \\
\addlinespace
\texttt{<head>} & La partie \emph{invisible} : informations pour le navigateur (encodage, titre de l'onglet, auteur). Les étudiants qui lisent ne voient jamais cela à l'écran. \\
\addlinespace
\texttt{<meta charset="UTF-8">} & Permet d'utiliser les accents et les caractères wolof (é, à, ñ\ldots). \textbf{Obligatoire pour le français.} \\
\addlinespace
\texttt{<title>} & Le texte affiché dans l'onglet du navigateur. \\
\addlinespace
\texttt{<body>} & La partie \emph{visible} : tout ce que l'utilisateur va voir sur la page. \\
\bottomrule
\end{tabular}
\end{center}

\begin{retenirbox}[À retenir]
Tout ce qui est \textbf{dans \texttt{<head>}} est invisible : c'est pour
le navigateur.\\
Tout ce qui est \textbf{dans \texttt{<body>}} est visible : c'est pour
l'utilisateur.
\end{retenirbox}

\subsection{Une balise, c'est quoi exactement ?}

Une \textbf{balise} (\emph{tag} en anglais) est un mot entouré de chevrons
\texttt{<\ldots>}. La plupart des balises marchent \textbf{par paire} : une
balise ouvrante et une balise fermante avec un \texttt{/}.

\begin{lstlisting}[language=HTML5]
<h1>Un titre</h1>
 |    |      |
 |    |      +-- balise fermante
 |    +-------- contenu
 +------------- balise ouvrante
\end{lstlisting}

Quelques balises sont \textbf{auto-fermantes} (elles n'ont pas besoin d'une
balise fermante car elles n'ont pas de contenu à entourer) :
\texttt{<img>}, \texttt{<br>}, \texttt{<meta>}, \texttt{<input>}.

\newpage

% ===========================================================================
\section{Les balises de contenu}
% ===========================================================================

HTML propose \textbf{plus de 100 balises}. Heureusement, on en utilise
régulièrement seulement une quinzaine. Voici les plus courantes --- celles
qu'on voit sur toute page web.

\subsection{Les titres}

HTML propose 6 niveaux de titres, de \texttt{<h1>} (le plus grand) à
\texttt{<h6>} (le plus petit).

\begin{lstlisting}[language=HTML5]
<h1>Université Cheikh Anta Diop</h1>
<h2>Faculté des Sciences</h2>
<h3>Département Informatique</h3>
<h4>DUT 1ère année</h4>
\end{lstlisting}

\begin{attentionbox}[Règle importante]
Il ne doit y avoir \textbf{qu'un seul \texttt{<h1>} par page} (le titre
principal). Les \texttt{<h2>}, \texttt{<h3>} peuvent apparaître plusieurs
fois pour structurer les sections et sous-sections.
\end{attentionbox}

\subsection{Les paragraphes}

\begin{lstlisting}[language=HTML5]
<p>L'École Supérieure Polytechnique forme des ingénieurs depuis 1964. Elle est située sur le campus de l'UCAD à Fann, Dakar.</p>
\end{lstlisting}

Chaque \texttt{<p>} représente \textbf{un paragraphe}. Le navigateur ajoute
automatiquement un espace entre deux paragraphes consécutifs.

\subsection{La mise en forme inline}

Pour mettre en avant un mot à l'intérieur d'un paragraphe :

\begin{lstlisting}[language=HTML5]
<p>Le cours est <strong>obligatoire</strong>. Les absences sont <em>vivement déconseillées</em>.</p>
\end{lstlisting}

\begin{itemize}[nosep]
    \item \texttt{<strong>} : texte \textbf{important} (affiché en gras par défaut)
    \item \texttt{<em>} : texte \emph{mis en valeur} (affiché en italique par défaut)
\end{itemize}

\subsection{Les listes}

\textbf{Listes non ordonnées} (à puces) avec \texttt{<ul>} :

\begin{lstlisting}[language=HTML5]
<ul>
    <li>Algorithmique</li>
    <li>Mathématiques</li>
    <li>Anglais</li>
</ul>
\end{lstlisting}

\textbf{Listes ordonnées} (numérotées) avec \texttt{<ol>} :

\begin{lstlisting}[language=HTML5]
<ol>
    <li>Installer VS Code</li>
    <li>Installer Git</li>
    <li>Créer un compte GitHub</li>
</ol>
\end{lstlisting}

Dans les deux cas, chaque élément est un \texttt{<li>} (\emph{list item}).

\subsubsection*{Listes imbriquées}

Une liste peut contenir une autre liste. Il suffit de mettre \texttt{<ul>}
ou \texttt{<ol>} \emph{à l'intérieur} d'un \texttt{<li>} :

\begin{lstlisting}[language=HTML5]
<ul>
    <li>Semestre 1
        <ul>
            <li>Algorithmique</li>
            <li>Mathématiques</li>
        </ul>
    </li>
    <li>Semestre 2
        <ul>
            <li>HTML / CSS / JavaScript</li>
            <li>Bases de données</li>
        </ul>
    </li>
</ul>
\end{lstlisting}

\begin{infobox}[L'indentation]
L'indentation (les espaces en début de ligne) n'est \textbf{pas obligatoire}
pour le navigateur, mais elle est \textbf{obligatoire pour vous}. Sans
indentation, un code HTML devient très vite illisible, même pour son auteur.
Prenez l'habitude dès maintenant.
\end{infobox}

\subsection{Les tableaux}

Un tableau HTML se construit ligne par ligne. Les balises essentielles :

\begin{itemize}[nosep]
    \item \texttt{<table>} : le tableau
    \item \texttt{<tr>} : une ligne (\emph{table row})
    \item \texttt{<th>} : une cellule d'en-tête (\emph{table header})
    \item \texttt{<td>} : une cellule de données (\emph{table data})
\end{itemize}

\begin{lstlisting}[language=HTML5]
<table>
    <tr>
        <th>Prénom</th>
        <th>Note</th>
    </tr>
    <tr>
        <td>Fatou</td>
        <td>15</td>
    </tr>
    <tr>
        <td>Moussa</td>
        <td>12</td>
    </tr>
</table>
\end{lstlisting}

\subsubsection*{Fusionner des cellules}

Pour fusionner plusieurs colonnes ou plusieurs lignes, on utilise les
attributs \texttt{colspan} (fusion horizontale) et \texttt{rowspan} (fusion
verticale) :

\begin{lstlisting}[language=HTML5]
<table>
    <tr>
        <th colspan="2">Lundi</th>  <!-- fusionne 2 colonnes -->
    </tr>
    <tr>
        <td>Algo</td>
        <td>HTML</td>
    </tr>
</table>
\end{lstlisting}

On reverra les tableaux en détail au Lab 2.

\subsection{Les liens}

Un lien vers une autre page se crée avec la balise \texttt{<a>}
(\emph{anchor}) et l'attribut \texttt{href} :

\begin{lstlisting}[language=HTML5]
<a href="https://ucad.sn">Site de l'UCAD</a>
<a href="https://esp.sn">Site de l'ESP</a>
\end{lstlisting}

Le texte entre \texttt{<a>} et \texttt{</a>} est ce que l'utilisateur voit.
L'URL dans \texttt{href} est la destination du clic.

\subsection{Les images}

\begin{lstlisting}[language=HTML5]
<img src="drapeau-senegal.png" alt="Drapeau du Sénégal">
\end{lstlisting}

\begin{itemize}[nosep]
    \item \texttt{src} : le chemin vers le fichier image (relatif ou URL complète)
    \item \texttt{alt} : description textuelle de l'image, \textbf{obligatoire}. Elle est lue par les personnes malvoyantes qui utilisent un lecteur d'écran, et affichée si l'image ne se charge pas.
\end{itemize}

\begin{attentionbox}[L'attribut \texttt{alt} n'est pas optionnel]
Une image sans \texttt{alt} est une image qui exclut les personnes
aveugles. Prenez l'habitude d'en mettre un à chaque \texttt{<img>}. C'est
un geste d'accessibilité --- pour certains étudiants malvoyants de votre
promotion, c'est la seule façon de «\,voir\,» votre page.
\end{attentionbox}

\newpage

% ===========================================================================
\section{Ce qu'on va construire : Jàng Afrig}
% ===========================================================================

Avant de plonger dans le code, prenons un moment pour regarder la montagne
qu'on va gravir ensemble.

\subsection{Le projet fil rouge}

\textbf{Jàng Afrig} (\textit{«\,apprendre l'Afrique\,»} en wolof) est une
application web de \textbf{quiz géographique} sur les 54 pays africains.
L'utilisateur :

\begin{enumerate}[nosep]
    \item Choisit un mode de jeu (capitales, drapeaux, langues, voisins\ldots)
    \item Répond à 10 questions à choix multiples
    \item Obtient un score à la fin
    \item Garde un classement local de ses meilleures performances
\end{enumerate}

\subsection{Les 5 pages de l'application}

\begin{center}
\begin{tabular}{@{}lp{9cm}@{}}
\toprule
\textbf{Page} & \textbf{Contenu} \\
\midrule
Accueil      & Choix du mode + saisie du pseudo + meilleur score \\
Quiz         & Énoncé d'une question + 4 propositions + barre de progression \\
Résultat     & Score final + bouton rejouer \\
Mes Scores   & Classement local top 5 par mode \\
À propos     & Sources des données, technologies utilisées \\
\bottomrule
\end{tabular}
\end{center}

\subsection{L'aventure par étapes}

On construit l'app \textbf{brique par brique}, à chaque séance :

\begin{center}
\begin{tabular}{@{}ll@{}}
\toprule
\textbf{Bloc de séances} & \textbf{Ce qu'on ajoute à Jàng Afrig} \\
\midrule
S1--S3 (HTML)    & Les 5 pages en HTML brut, sans style \\
S4--S6 (CSS)     & Design panafricain, responsive mobile \\
S7--S10 (JS)     & Premier quiz interactif avec 10 pays hardcodés \\
S11--S13 (API)   & Connexion à REST Countries, 54 pays réels, 4 modes \\
S14 (polish)     & Chrono, \texttt{localStorage}, finitions \\
S15 (Hybride)    & Déploiement en ligne + démo collective \\
\bottomrule
\end{tabular}
\end{center}

\begin{projetbox}[À la fin du semestre]
Chacun de vous aura son application accessible sur Internet via une URL
publique du type \texttt{https://\emph{votre-username}.github.io/jang-afrig/}.
Vous pourrez la partager à votre famille par WhatsApp.
\end{projetbox}

\newpage

% ===========================================================================
\section{Le workflow du cours}
% ===========================================================================

\subsection{Le rythme hebdomadaire}

Chaque séance suit toujours la même mécanique :

\begin{center}
\begin{tabular}{@{}lp{9cm}@{}}
\toprule
\textbf{Type de séance} & \textbf{Déroulement} \\
\midrule
\textbf{CM} & Rappel $\to$ analogie $\to$ notions nouvelles avec démo du prof $\to$ mini-exercice en classe $\to$ teaser du Lab suivant \\
\addlinespace
\textbf{Lab} & Rappel du CM $\to$ démo de l'objectif du jour $\to$ vous codez en pas-à-pas $\to$ exercices d'application $\to$ \textbf{routine Git} $\to$ annonce de la séance suivante \\
\bottomrule
\end{tabular}
\end{center}

\subsection{La règle d'or : 1 Lab = 1 tag Git}

À la fin de chaque Lab, vous poserez un \textbf{tag Git} qui marque une
version de votre application. Au bout de 10 Labs, votre historique Git
ressemblera à ceci :

\begin{lstlisting}
v0     Setup initial (Lab 0)
v0.1   Première page HTML (Lab 1)
v0.2   HTML structurant (Lab 2)
v0.3   CSS de base (Lab 3)
...
v0.9   Tous les modes avec API (Lab 9)
v0.14  Polish final (Lab 10)
v1.0   Déploiement en ligne (S15)
\end{lstlisting}

\begin{retenirbox}[Chaque Lab = une brique posée]
Si vous venez à manquer une séance, vous aurez \textbf{un repo en retard
d'une brique}. Facile à combler en relisant le PDF du Lab manqué. Mais il
faut rattraper \textbf{avant} la séance suivante, sinon la dette s'accumule.
\end{retenirbox}

\subsection{Le rapport à la fin de chaque Lab}

Vous avez produit votre premier rapport pour le Lab 0. Même routine pour
les 10 Labs à venir :

\begin{itemize}[leftmargin=*,nosep]
    \item 1 à 2 pages maximum
    \item 4 sections : \textbf{ce que j'ai fait}, \textbf{difficultés}, \textbf{comment j'ai résolu}, \textbf{lien repo}
    \item Déposé dans \texttt{rapports/rapport-lab-X-NOM-Prenom.pdf}
\end{itemize}

C'est votre \textbf{carnet de bord} : un petit effort à chaque Lab, un vrai
parcours à la fin du semestre.

\newpage

% ===========================================================================
\section{Routine Git : rappel}
% ===========================================================================

Vous avez appris au Lab 0 les 5 étapes du cycle Git. Rappel rapide --- on
va la répéter 10 fois dans le semestre.

\begin{center}
\begin{tabular}{@{}rl@{}}
\toprule
\textbf{Étape} & \textbf{Commande} \\
\midrule
1. Modifier des fichiers dans VS Code & \emph{(on édite le code)} \\
2. Voir ce qui a changé                & \texttt{git status} \\
3. Préparer tous les changements       & \texttt{git add .} \\
4. Enregistrer avec un message         & \texttt{git commit -m "v0.X : titre"} \\
5. Envoyer sur GitHub                  & \texttt{git push} \\
6. Poser un tag de version             & \texttt{git tag v0.X} \\
7. Envoyer le tag                      & \texttt{git push --tags} \\
\bottomrule
\end{tabular}
\end{center}

\begin{retenirbox}[Mémoire musculaire]
Au bout de 3 ou 4 Labs, vous taperez cette séquence sans réfléchir. C'est
le but. Git doit devenir un \textbf{réflexe}, pas un obstacle.
\end{retenirbox}

% ===========================================================================
\section{Mini-exercice en classe}
% ===========================================================================

\begin{objectifbox}[Exercice guidé --- écrit ensemble]
On va écrire \textbf{ensemble}, ici et maintenant, notre première page HTML.

\medskip

\textbf{Énoncé.} Créer une page \texttt{moi.html} qui contient :
\begin{enumerate}[nosep]
    \item Un titre principal \texttt{<h1>} avec votre prénom
    \item Un paragraphe de présentation avec un \texttt{<strong>} quelque part
    \item Une liste \texttt{<ul>} de vos 3 matières préférées
    \item Un lien vers le site de l'UCAD
\end{enumerate}
\end{objectifbox}

\subsection{Solution attendue}

\begin{lstlisting}[language=HTML5]
<!DOCTYPE html>
<html lang="fr">
    <head>
        <meta charset="UTF-8">
        <title>Ma page</title>
    </head>
    <body>
        <h1>Fatou</h1>
        <p>Je suis étudiante en <strong>DUT1-Informatique</strong> à l'ESP, UCAD.</p>
        <ul>
            <li>Algorithmique</li>
            <li>Mathématiques</li>
            <li>HTML/CSS/JS</li>
        </ul>
        <a href="https://ucad.sn">Site de l'UCAD</a>
    </body>
</html>
\end{lstlisting}

\begin{infobox}[Levez la main]
Si vous n'arrivez pas à l'écrire en 10 minutes, \textbf{levez la main}.
C'est normal : c'est votre première page. On la refera calmement ensemble
au Lab 1 (séance prochaine).
\end{infobox}

% ===========================================================================
\section{Résumé et transition}
% ===========================================================================

\subsection{Ce qu'on a vu}

\begin{itemize}[leftmargin=*,nosep]
    \item Une page web = \textbf{HTML} (structure) + \textbf{CSS} (déco) + \textbf{JS} (comportement)
    \item Un \textbf{navigateur} demande une page à un \textbf{serveur} qui répond avec du HTML
    \item Une page HTML a un \texttt{<head>} (invisible) et un \texttt{<body>} (visible)
    \item Les balises essentielles : \texttt{h1}-\texttt{h6}, \texttt{p}, \texttt{strong}, \texttt{em}, \texttt{ul}/\texttt{ol}/\texttt{li}, \texttt{table}/\texttt{tr}/\texttt{th}/\texttt{td}, \texttt{a}, \texttt{img}
    \item Jàng Afrig se construira \textbf{brique par brique} pendant 15 séances
    \item 1 Lab = 1 commit + 1 tag Git + 1 rapport
\end{itemize}

\subsection{Ce qu'on fera au Lab 1}

\begin{itemize}[leftmargin=*,nosep]
    \item Créer \texttt{index.html} dans votre repo \texttt{jang-afrig}
    \item Construire la \textbf{page d'accueil} de Jàng Afrig en HTML brut
    \item Travailler les listes imbriquées
    \item Travailler les tableaux avec \texttt{colspan}/\texttt{rowspan}
    \item Premier tag Git : \texttt{v0.1}
\end{itemize}

\subsection{Consigne pour la prochaine séance}

\begin{attentionbox}[Avant le Lab 1]
\begin{enumerate}[nosep]
    \item \textbf{Lisez} le PDF du Lab 1 que le prof vous a distribué (15 min de lecture).
    \item \textbf{Vérifiez} que votre environnement fonctionne toujours (\texttt{code .} dans votre dossier \texttt{jang-afrig}, Live Server opérationnel).
    \item \textbf{Prévoyez} une capture d'écran d'un drapeau africain de votre choix pour les exercices.
\end{enumerate}
Aucun devoir de code. Le Lab 1 démarrera du vide (page d'accueil à partir de zéro).
\end{attentionbox}

\vspace{1cm}

\begin{center}
\fbox{
\begin{minipage}{0.85\textwidth}
\centering
\vspace{0.5cm}
\textbf{\Large CM1 --- Terminé \cmark}\\[0.3cm]
\normalsize
Vous savez comment une page web fonctionne, et comment l'écrire.\\
Prochaine étape : \textbf{la construire}.\\[0.2cm]
\textit{Rendez-vous au Lab 1 !}\\[0.3cm]
\vspace{0.3cm}
\end{minipage}
}
\end{center}

\end{document}
